\documentclass[11pt]{article}
\usepackage[a4paper,margin=1in]{geometry}
\usepackage{amsmath,amssymb,amsthm,mathtools,booktabs}
\usepackage[T1]{fontenc}
\usepackage{lmodern,microtype}
\usepackage[hidelinks]{hyperref}
\newtheorem{theorem}{Theorem}
\newtheorem{lemma}[theorem]{Lemma}
\newtheorem{proposition}[theorem]{Proposition}
\newtheorem{remark}[theorem]{Remark}
\newcommand{\R}{\mathbb R}
\DeclareMathOperator{\rank}{rank}
\title{A strict pointwise tensor-train approximation bound\newline by finite singular-subspace tie handling}
\author{Matthew J. Colbrook\thanks{Department of Applied Mathematics and Theoretical Physics, University of Cambridge, Cambridge, United Kingdom. Email: m.colbrook@damtp.cam.ac.uk.}}
\date{11 September 2026}
\hypersetup{pdfauthor={Matthew J. Colbrook}}
\usepackage{xurl}
\setlength{\emergencystretch}{3em}
\begin{document}
\maketitle
\begin{center}\small Recovered qualified manuscript for TR-04. Not a uniform-factor improvement.\end{center}
\begin{quote}\small Authorship recorded at the submitter\textquotesingle s request. This recovered draft was generated with AI assistance. The complete original source passed independent agent review against the exact repository target. This is not external human peer review or formal certification; no priority claim is made. Original provenance, full-source review hashes and diagnostic results accompany this manuscript in the submission record.\end{quote}
\begin{abstract}
For prescribed tensor-train rank bounds, a deterministic algorithm in
the exact arithmetic/SVD model returns an approximation with squared
error strictly smaller than $(d-1)$ times the optimal squared error
whenever that optimum is positive, and reconstructs exactly at zero
optimum. It performs at most $n_1$ tensor-train SVD completions. The only
modification to the standard construction is to test a finite collection
of equally optimal first-cut singular subspaces when a positive singular
value straddles the truncation boundary. This addresses the strict
pointwise inequality displayed in TR-04. It does not produce a uniform
constant $c<d-1$: an explicit fixed-format family shows that its error
ratios can approach $d-1$ already when $d=3$.
\end{abstract}

\section{The target and the distinction between two guarantees}
Let $d\ge3$, let $A\in\R^{n_1\times\cdots\times n_d}$ be a dense tensor,
and let $r_1,\ldots,r_{d-1}$ be positive integers. Denote by $S_r$ the
set of tensors whose unfolding at every cut $1,\ldots,j\mid j+1,\ldots,d$
has rank at most $r_j$. Write
\[
 E_* = E_*(A,r)=\min_{Y\in S_r}\|A-Y\|_F^2.
\]
This minimum exists: $S_r$ is a closed intersection of determinantal
sets, and a minimizing sequence may be confined to a bounded ball.
The displayed strict inequality in the repository is
\begin{equation}\label{eq:target}
 \|A-X\|_F^2<(d-1)E_*\qquad(E_*>0),
\end{equation}
with exact reconstruction when $E_*=0$, in the idealized arithmetic/SVD
model \cite{repo,workshop}.
A guarantee of the form \eqref{eq:target} for every input does not imply
that the supremum of all error ratios is smaller than $d-1$. We prove
the former, and make no claim of the latter. In particular this note
must not be presented as a smaller uniform worst-case constant.

\section{A completion estimate}
The standard TT-SVD guarantee \cite[Theorem 2.2 and Corollary 2.4]{O}
will be used on the remaining cuts after a first projection.

\begin{lemma}\label{lem:completion}
Let $U\in\R^{n_1\times k}$ have orthonormal columns, with $k\le r_1$,
and put $P=UU^T$. Compress the first mode to form $B=U^T A$.
Treat its first two indices as a single index, so $B$ is a tensor of
order $d-1$ and format $(kn_2)\times n_3\times\cdots\times n_d$.
Apply ordinary TT-SVD to this tensor with cut-rank bounds
$r_2,\ldots,r_{d-1}$, and lift the result through $U$ to obtain $X_U$.
Then $X_U\in S_r$, and for every $Y\in S_r$,
\begin{equation}\label{eq:completion}
 \|A-X_U\|_F^2
 \le \|(I-P)A\|_F^2+(d-2)\|P(A-Y)\|_F^2.
\end{equation}
Here and below a matrix acting on a tensor acts on its first mode.
\end{lemma}
\begin{proof}
The first cut of $X_U$ has rank at most $k$. The remaining cuts have the
specified rank bounds, which are preserved by lifting through $U$.
The compressed tensor $U^TY$ satisfies every remaining cut constraint:
compression is left multiplication on each such unfolding and cannot
increase its rank. The TT-SVD guarantee on the $d-2$ remaining cuts is
therefore
\[
 \|B-\widehat B\|_F^2\le(d-2)\|U^T(A-Y)\|_F^2.
\]
For completeness, this usual guarantee follows by summing the squared
orthogonal truncation residuals at successive cuts. Each residual is at
most the best error at that cut; projecting any feasible comparison
tensor through the previously computed orthogonal factors shows that
each is at most the squared error of that comparison tensor. There are
$d-2$ cuts. This argument also gives zero residuals when the input is
already feasible.
Finally, the first-mode discarded component is orthogonal to the lifted
completion residual, giving
\[
 \|A-X_U\|_F^2=\|(I-P)A\|_F^2+\|B-\widehat B\|_F^2.
\]
The identity $\|U^T(A-Y)\|_F=\|P(A-Y)\|_F$ completes the proof.
\end{proof}

\section{The finite algorithm}
Let $A_{(1)}$ denote the $n_1\times(n_2\cdots n_d)$ first unfolding.
If $A=0$, return $0$. Compute its SVD and let
\[
 \varepsilon^2=\sum_{i>r_1}\sigma_i(A_{(1)})^2,
\]
where singular values beyond the matrix dimensions are interpreted as
zero. Construct a list $\mathcal P$ of first-mode projectors as follows.

If $\rank A_{(1)}\le r_1$, use just the projector onto the column space
of $A_{(1)}$. Otherwise let $\tau=\sigma_{r_1}(A_{(1)})>0$. Let $H$ be
the span of left singular vectors with singular value strictly greater
than $\tau$, put $h=\dim H$, and let $E$ be the left singular space for
$\tau$, of dimension $t$. Set $s=r_1-h$, so $1\le s\le t$.
If $s=t$, use just the projector onto $H\oplus E$.
If $s<t$, take any fixed orthonormal basis $u_0,\ldots,u_{t-1}$ of $E$
and form the $t$ cyclic windows
\[
 W_j=\operatorname{span}\{u_j,u_{j+1},\ldots,u_{j+s-1}\},
 \qquad 0\le j<t,
\]
with subscripts modulo $t$. Use the $t$ projectors
\begin{equation}\label{eq:projectors}
 P_j=P_H+P_{W_j}.
\end{equation}
Each projector in $\mathcal P$ is a best first-cut rank-$r_1$ projection
(or an exact column-space projection when the rank is smaller).
For each $P\in\mathcal P$, choose an orthonormal column matrix $U$ with
$P=UU^T$, form $X_U$ by Lemma~\ref{lem:completion}, and compute its
actual Frobenius error. Return a candidate of minimum error, resolving
any final tie deterministically.

There are at most $n_1$ candidates. SVDs, matrix products, tensor
reshapings, and error computations have polynomial cost in the dense
input size and the rank parameters. Thus this is a deterministic
polynomial-time algorithm in the arithmetic/SVD model. Equality of
singular values is tested exactly in that model. A floating-point
implementation with a tolerance is an illustration, not a bit-complexity
implementation of the exact tie test.

\section{Proof of the strict guarantee}
We first isolate the only equality case requiring attention.

\begin{lemma}\label{lem:matrix}
Suppose $\rank A_{(1)}>r_1$, and let $Y$ satisfy
$\rank Y_{(1)}\le r_1$ and
$\|A-Y\|_F^2=\varepsilon^2$.
With $H$ and $E$ defined above, the residual $R=A-Y$ obeys
$P_HR=0$. If $s=t$, it also obeys $(P_H+P_E)R=0$.
\end{lemma}
\begin{proof}
Let $S$ be the column space of $Y_{(1)}$. The orthogonal decomposition
\[
 \|A_{(1)}-Y_{(1)}\|_F^2
 =\|(I-P_S)A_{(1)}\|_F^2+
   \|P_SA_{(1)}-Y_{(1)}\|_F^2
\]
shows that equality with the best rank-$r_1$ error forces
$Y_{(1)}=P_SA_{(1)}$ and forces $P_S$ to maximize the captured singular
energy. Since $\tau>0$, $\dim S=r_1$. The equality characterization of
this variational maximum gives
$S=H\oplus E_0$ for an $s$-dimensional subspace $E_0\subset E$.
One can see this directly by writing
$\operatorname{tr}(P_SA_{(1)}A_{(1)}^T)$ in an eigenbasis: all weights
at eigenvalues strictly above $\tau^2$ must be one, all below must be
zero, and the remaining weight is $s$ in the $\tau^2$ eigenspace.
Thus $H\subset S$, and also $E\subset S$ when $s=t$.
The asserted residual orthogonality follows.
\end{proof}

\begin{theorem}\label{thm:strict}
The finite algorithm above returns $X\in S_r$ satisfying
\eqref{eq:target} whenever $E_*>0$, and returns $X=A$ when $E_*=0$.
\end{theorem}
\begin{proof}
Fix a best feasible $Y\in S_r$, and write $R=A-Y$.
The first-cut lower bound is $\varepsilon^2\le E_*$.
Every tested projector $P$ has $\|(I-P)A\|_F^2=\varepsilon^2$;
Lemma~\ref{lem:completion} gives
\begin{equation}\label{eq:basic}
 \|A-X_U\|_F^2
 \le\varepsilon^2+(d-2)\|PR\|_F^2
 \le\varepsilon^2+(d-2)E_*.
\end{equation}
If $\varepsilon^2<E_*$, this is strictly smaller than $(d-1)E_*$.
In particular, this covers $\rank A_{(1)}\le r_1$ whenever $E_*>0$.

It remains to consider $\varepsilon^2=E_*>0$. By
Lemma~\ref{lem:matrix}, $P_HR=0$. If $s=t$, the single tested
projector annihilates $R$, so \eqref{eq:basic} bounds the error by
$E_*<(d-1)E_*$. If $0<s<t$, the cyclic-window projectors satisfy the
exact averaging identity
\[
 \frac1t\sum_{j=0}^{t-1}P_{W_j}=\frac{s}{t}P_E,
\]
since each basis vector appears in exactly $s$ windows. Hence
\[
 \frac1t\sum_{j=0}^{t-1}\|P_jR\|_F^2
 =\frac{s}{t}\|P_ER\|_F^2
 \le\frac{s}{t}E_*<E_*.
\]
At least one tested projector has this bound. Its completion satisfies
\[
 \|A-X_U\|_F^2
 \le\left(1+(d-2)\frac{s}{t}\right)E_*
 <(d-1)E_*.
\]
The returned minimum-error candidate is no worse.

Finally, if $E_*=0$, then $A\in S_r$, $\varepsilon^2=0$, and
\eqref{eq:basic} with $R=0$ forces the returned error to be zero.
\end{proof}

\section{Why this is not a uniform constant improvement}
This distinction is material to the scope of the result, not a numerical
precision issue. Consider $d=3$, format $3\times3\times2$, ranks $(2,1)$,
and the tensor
\[
 A=\alpha e_1\otimes e_1\otimes f_1
   +\beta e_2\otimes e_2\otimes f_1
   +\gamma e_3\otimes e_3\otimes f_2,
 \qquad \alpha>\gamma\ge\beta>0.
\]
Its optimal error is $E_*=\gamma^2$.
Indeed, a tensor feasible for ranks $(2,1)$ is $M\otimes v$ with
$\rank M\le2$. For a unit vector $v=(c,s)$, the contracted matrix is
$\operatorname{diag}(\alpha c,\beta c,\gamma s)$. The maximum possible
sum of its two largest squared singular values, also maximizing over
$c^2+s^2=1$, is
\[
 \max\{\alpha^2+\beta^2,\alpha^2,\beta^2,\gamma^2\}
 =\alpha^2+\beta^2.
\]
This proves the claimed optimum by orthogonal projection.

When $\gamma>\beta$, the best first-cut rank-two subspace is uniquely
$\operatorname{span}(e_1,e_3)$. The remaining rank-one truncation retains
only the $\alpha$ term. The algorithm's squared error is therefore
$\beta^2+\gamma^2$, and its ratio to optimum is
\[
 1+(\beta/\gamma)^2<2.
\]
Keeping $\alpha=2$, $\beta=1$ and sending $\gamma\downarrow1$ shows that
these ratios approach $2=d-1$ in a fixed format. At $\gamma=\beta$,
an unfortunate ordinary TT-SVD tie choice can attain ratio exactly two;
the finite tie handling removes that equality. No theorem here asserts
a uniform ratio bounded by $2-\delta$ for a fixed $\delta>0$.

\section{Reproducibility and scope}
The accompanying \texttt{verify\_tr04.py} checks the projector averaging
identity exactly, checks the fixed-format family and rotated tie bases,
and tests the numerical implementation for rank feasibility and exact
reconstruction up to floating-point error. The exact mathematical
algorithm and its guarantee are Theorem~\ref{thm:strict}; a tolerance-based
SVD is not claimed to detect exact ties in finite precision.
No ranks are enlarged. No oracle for $E_*$ or for an optimal $Y$ is used:
those objects appear only in the proof. The assertion addressed is the
strict pointwise target, not a stronger uniform-factor or bit-complexity
formulation.

\begin{thebibliography}{9}
\bibitem{O}
I.~V.~Oseledets, \emph{Tensor-Train Decomposition},
SIAM Journal on Scientific Computing 33 (2011), 2295--2317,
Theorem 2.2 and Corollary 2.4.
\href{https://doi.org/10.1137/090752286}{doi:10.1137/090752286}.
\bibitem{workshop}
N.~Amsel et al., \emph{Linear Systems and Eigenvalue Problems:
Open Questions from a Simons Workshop},
\href{https://arxiv.org/abs/2602.05394}{arXiv:2602.05394v3} (2026),
Section 6.1, Problem 6.1.
\bibitem{repo}
OpenProblemsInNLA, \emph{TR-04: Improve the worst-case approximation
factor for prescribed tensor-train ranks}, canonical statement.
\url{https://github.com/ajt60gaibb/OpenProblemsInNLA/tree/main/tensor-computations/TR-04}.
\end{thebibliography}
\end{document}
